% !TEX root = ../main.tex

% 专硕请取消下面注释
\makeatletter
  \def\hitsz@csubjecttitle{类别}
\makeatother

\hitszsetup{
  %******************************
  % 注意：
  %   1. 配置里面不要出现空行
  %   2. 不需要的配置信息可以删除
  %******************************
  %
  %=====
  % 秘级
  %=====
  statesecrets={公开},
  natclassifiedindex={TM301.2},
  intclassifiedindex={62-5},
  %
  %=========
  % 中文信息
  %=========
  ctitleone={基于图神经网络的多行为},%本科生封面使用
  ctitletwo={去噪推荐方法研究},%本科生封面使用
  ctitlecover={基于图神经网络的多行为去噪推荐方法研究},%放在封面中使用，自由断行
  ctitle={基于图神经网络的多行为去噪推荐方法研究},%放在原创性声明中使用
  csubtitle={一条副标题}, %一般情况没有，可以注释掉
  cxueke={专业},
  cpostgraduatetype={专业},
  csubject={计算机技术},
  % csubject={机械工程},
  % caffil={计算机科学与技术学院},
  caffil={哈尔滨工业大学（深圳）},
  cauthor={樊红雨},
  csupervisor={马江虹副教授},
  % cassosupervisor={某某某 教授}, % 副导师
  % ccosupervisor={某某某 教授}, % 合作导师
  % 日期自动使用当前时间，若需指定按如下方式修改：
  cdate={2026年1月},
  % 指定第二页封面的日期，即答辩日期
  cdatesecond={2026年01月},
  cstudentid={23S151008},
  cstudenttype={同等学力人员}, %非全日制教育申请学位者
  %（同等学力人员）、（工程硕士）、（工商管理硕士）、
  %（高级管理人员工商管理硕士）、（公共管理硕士）、（中职教师）、（高校教师）等
  %
  %
  %=========
  % 英文信息
  %=========
  etitle={Research on Multi-Behavior Denoising Recommendation Method Based on Graph Neural Networks},
  esubtitle={This is the sub title},
  exueke={Engineering},
  esubject={Electronic and Information Engineering},
  eaffil={Harbin Institute of Technology, Shenzhen},
  eauthor={Hongyu Fan},
  esupervisor={Associate Prof. Jianghong Ma},
  % eassosupervisor={XXX},
  % ecosupervisor={Prof. XXX}, % Co-Supervisor off Campus
  % 日期自动生成，若需指定按如下方式修改：
  edate={January, 2026},
  estudenttype={Master of Engineering},
  %
  % 关键词用“英文逗号”分割
  ckeywords={推荐系统, 图神经网络, 去噪推荐, 多行为推荐},
  ekeywords={recommender system, graph neural network, denoising recommendation, multi-behavior recommendation},
}

% 中文摘要
\begin{cabstract}

随着用户交互行为类型的日益丰富，推荐系统在多行为场景下面临严重的噪声干扰与偏好混叠问题。用户在电商或内容平台上的点击、收藏、加购、购买等行为虽然能够反映潜在兴趣，但其中包含的随机、误触或弱相关交互会削弱模型对真实偏好的刻画能力。此外，传统多行为推荐多依赖协同信号，缺乏模态语义信息的辅助，难以在复杂语义环境下准确识别噪声行为。针对上述问题，本文围绕“多行为协同去噪”与“多模态语义引导去噪”两个方向展开研究，提出了两个具有互补性的去噪推荐模型。

首先，提出噪声识别与偏好提取的多行为去噪推荐方法，面向多行为推荐中的噪声干扰问题。该模型从行为间与行为内两个层面建模用户偏好关系，设计了行为加权机制与自监督约束，从多行为图结构中自动识别并抑制噪声交互的影响。同时，引入噪声感知的偏好提取模块，使模型在保留高质量辅助行为信号的同时，消除弱相关或不可靠的交互，从而获得更精确的偏好表示。

其次，为进一步扩展至多模态多行为场景，提出基于噪声识别与偏好提取的多模态多行为去噪推荐方法。该方法引入大型视觉语言模型提取物品的图像与文本语义表示，通过模态语义相似度实现硬去噪，去除语义不一致的辅助交互；在图神经网络框架中，设计语义感知注意力机制以实现软去噪，自适应调整邻居节点的信息权重，从而聚焦于语义一致的高质量交互。此外，提出跨图对比学习模块，在语义增强图与去噪图之间建立一致性约束，实现模态语义信息向结构空间的有效迁移。

本文的创新性体现在结合协同信号建模与多模态语义理解技术，提出两种创新方法以解决多行为推荐场景中的噪声干扰与真实偏好识别问题。
一方面，从协同信息建模与去噪机制两方面入手，设计了基于行为内与行为间双层结构的噪声感知偏好提取模型，实现对真实偏好与噪声交互的有效区分，提升了多行为推荐的精确性与鲁棒性；
另一方面，利用大型视觉语言模型的语义理解能力与多模态特征融合机制，构建多模态多行为去噪推荐框架，通过语义一致性约束实现软、硬去噪的协同优化，从视觉与文本模态中挖掘跨模态真实偏好信号，为复杂交互场景下的鲁棒推荐提供了新的思路与方法。


\end{cabstract}

% 英文摘要
\begin{eabstract}

With the increasing diversity of user interaction behaviors, recommender systems in multi-behavior scenarios face significant challenges of noise interference and preference entanglement. User actions on e-commerce or content platforms—such as clicks, favorites, add-to-cart, and purchases—can reflect potential interests, yet they often include random, accidental, or weakly related interactions that obscure true user preferences. Moreover, traditional multi-behavior recommendation methods primarily rely on collaborative signals, lacking semantic guidance from multimodal information, which limits their ability to identify noisy behaviors in complex semantic environments.

To address these issues, this thesis investigates two complementary directions: collaborative denoising and multimodal semantic-guided denoising, and proposes two novel denoising recommendation models.

First, a Noise-aware Inter- and Intra-behavior Preference Extraction model (NIPER) is proposed to address noise interference in multi-behavior recommendation. This model captures user preference relationships at both inter-behavior and intra-behavior levels, and introduces a behavior-weighted mechanism together with a self-supervised constraint to automatically identify and suppress noisy interactions in the multi-behavior graph. Furthermore, a noise-aware preference extraction module is designed to preserve high-quality auxiliary behavioral signals while filtering out weakly relevant or unreliable interactions, thus yielding more accurate user preference representations.

Second, to further extend the denoising framework to multimodal multi-behavior scenarios, a large-model semantic-guided multimodal denoising recommendation framework is proposed. This method leverages large vision-language models to extract semantic representations from both textual and visual modalities of items. Based on cross-modal semantic similarity, a hard-denoising mechanism removes inconsistent auxiliary interactions. Within the graph neural network (GNN) architecture, a semantic-aware attention mechanism is introduced to achieve soft denoising by adaptively adjusting neighbor weights, thereby emphasizing semantically consistent and informative interactions. In addition, a cross-graph contrastive learning module establishes consistency constraints between the denoised graph and the semantic-augmented graph, enabling effective transfer of semantic information into the structural representation space.

The main contributions of this thesis are as follows:
(1) A collaborative denoising method that distinguishes true preferences from noisy interactions through a dual-level (inter- and intra-behavior) preference extraction framework, enhancing the accuracy and robustness of multi-behavior recommendation;
(2) A multimodal semantic-guided denoising framework that integrates large vision-language models and attention-based message passing, achieving synergistic optimization of hard and soft denoising via semantic consistency constraints. This design effectively captures cross-modal preference signals from textual and visual modalities, providing a new perspective and methodology for robust recommendation in complex multi-behavior and multimodal environments.


\end{eabstract}
